%% toc: P15 160 160.1,160.2,160.3,160.4
% Proofs for Feller's numerical record of Stirling's approximation (Feller I, 3rd ed., II.9, pp. 52--54)
% and for Durrett's Theorem 3.1.3 as a case of the i.i.d. central limit theorem (Durrett PTE5, pp. 113, 115);
% three convolution details (Durrett PTE5, Theorems 2.1.15, 2.1.18, 2.1.20, pp. 49--52).
\section{Stirling's approximation, the fair-sign normal limit and three convolution details}
\label{sec:annex-p15-w02c-p15-011-stirling-record}

\begin{annexitem}{annex:p15-w02c-p15-011-stirling-numerics}{thm:p15-w02c-p15-011-stirling-numerics}{Brackets and interval values}
\begin{proof}
\emph{Two elementary inequalities.} For real \(x>0\),
\[
 x-\frac{x^2}2<1-e^{-x}<x .
\]
Put \(g(x)=1-e^{-x}-x\). Then \(g(0)=0\) and \(g'(x)=e^{-x}-1<0\) for
\(x>0\), so \(g(x)<0\); this is the right inequality, and it also says
\(e^{-x}>1-x\). Put \(h(x)=1-e^{-x}-x+x^2/2\). Then \(h(0)=0\) and
\(h'(x)=e^{-x}-1+x>0\) for \(x>0\) by the inequality just proved, so
\(h(x)>0\); this is the left inequality.

\emph{Item~1 for \(n\geq2\).} By \cref{thm:p15-stirling-with-remainder},
\(A_n/n!=e^{-r_n}\) with \(1/(12n+1)<r_n<1/(12n)\). The function
\(x\mapsto1-e^{-x}\) is strictly increasing, so with
\(x_-=1/(12n+1)\) and \(x_+=1/(12n)\) the two elementary inequalities give
\[
 x_--\frac{x_-^2}2<1-e^{-x_-}<1-\frac{A_n}{n!}<1-e^{-x_+}<x_+ .
\]
The four cases are, as fractions (multiply by \(100\) for per cent):
\begin{itemize}
\item \(n=2\): \(x_-=1/25\), \(x_--x_-^2/2=1/25-1/1250=0.0392\), and
      \(x_+=1/24=0.041666\ldots<0.0417\);
\item \(n=5\): \(x_-=1/61\), \(x_--x_-^2/2=1/61-1/7442>0.016393-0.000135
      =0.016258>0.0162\), and \(x_+=1/60=0.016666\ldots<0.0167\);
\item \(n=10\): \(x_-=1/121\),
      \(x_--x_-^2/2=1/121-1/29282>0.0082644-0.0000342=0.0082302>0.00823\),
      and \(x_+=1/120=0.0083333\ldots<0.00834\);
\item \(n=100\): \(x_-=1/1201\), \(x_+=1/1200=0.00083333\ldots<0.000834\),
      and, since \(1/1201>0.00083263\) and \(1/2884802<0.00000035\),
      \[
       x_--\frac{x_-^2}2=\frac1{1201}-\frac1{2884802}
       >0.00083228>0.000832 .
      \]
\end{itemize}
This proves the stated intervals for \(n=2,5,10,100\).

\emph{Interval evaluation.} The closed forms
\[
 A_n=\sqrt{2\pi}\,n^{n+1/2}e^{-n},\qquad
 A_ne^{1/(12n)}=\sqrt{2\pi}\,n^{n+1/2}e^{-n+1/(12n)}
\]
are strictly increasing in \(\pi\) and, because \(-n<-n+1/(12n)<0\),
strictly decreasing in \(e\). The decimal bounds
\(3.141592653<\pi<3.141592654\) and \(2.718281828<e<2.718281829\)
therefore give lower bounds by evaluating at the smaller value of \(\pi\)
and the larger value of \(e\), and upper bounds the other way. Evaluating at
these rational endpoints gives
\[
 \begin{array}{c|c|c}
  n & A_n & A_ne^{1/(12n)}\\\hline
  1 & (0.9221370086,\ 0.9221370092) & \\
  2 & (1.9190043505,\ 1.9190043523) & (2.0006520467,\ 2.0006520485)\\
  5 & (118.01916782,\ 118.01916807) & (120.00263695,\ 120.00263720)\\
  10 & (3{,}598{,}695.611,\ 3{,}598{,}695.626) &
 \end{array}
\]
For \(n=1\), \(1!=1\) and the relative error lies in
\((0.0778629908,\ 0.0778629914)\), which is inside \((0.07786,0.07787)\);
this completes item~1. The intervals for \(A_1\), \(A_2\), \(A_5\) and
\(A_{10}\) give the leading digits displayed in item~2. The first three
lie in
\([0.92205,0.92215)\), \([1.9185,1.9195)\) and \([118.0185,118.0195)\)
respectively, so the rounded values are \(0.9221\), \(1.919\) and
\(118.019\). The \(n=10\) interval also gives
\(10!-A_{10}\in(30{,}104.374,\ 30{,}104.389)\), used in
\cref{def:p15-w02c-p15-011-relative-not-absolute}. The third column
proves item~4.

\emph{Rounding of Feller's percentages.} The certified intervals lie in the
rounding cells of Feller's figures: \((7.786,7.787)\subset[7.5,8.5)\),
\((3.92,4.17)\subset[3.5,4.5)\), \((1.62,1.67)\subset[1.5,2.5)\),
\((0.823,0.834)\subset[0.75,0.85)\) and
\((0.0832,0.0834)\subset[0.075,0.085)\). Each figure is therefore the
correctly rounded value at the precision printed.

\emph{Item~3.} Put \(x_n=1/(12n)-1/(12n+1)=1/(12n(12n+1))>0\). Since
\(n!=A_ne^{r_n}\),
\[
 \frac{A_ne^{1/(12n)}}{n!}=e^{1/(12n)-r_n},
 \qquad 0<\frac1{12n}-r_n<\frac1{12n}-\frac1{12n+1}=x_n,
\]
by the Robbins bounds. The exponential is strictly increasing, so
\(1<A_ne^{1/(12n)}/n!<e^{x_n}\); subtracting one gives the first two
inequalities. For the third, the mean value theorem gives
\(e^{x}-1=xe^{\xi}\) for some \(\xi\in(0,x)\), so \(e^x-1<xe^x\) for
\(x>0\). Also \(e^y<1/(1-y)\) for \(0<y<1\), because
\(e^{-y}>1-y\) by the first paragraph. Since \(x_n\leq x_1=1/156\),
\[
 e^{x_n}-1<x_ne^{x_n}<\frac{156}{155}\,x_n
   <\frac{156}{155}\cdot\frac1{144n^2}
   =\frac{156}{22320\,n^2}<\frac{0.007}{n^2},
\]
because \(12n(12n+1)>144n^2\) and \(0.007\cdot22320=156.24>156\).
In per cent this is less than \(0.70\,n^{-2}\).

The ratio of the lower member to the upper member is
\(e^{1/(12n+1)-1/(12n)}=e^{-x_n}\). By the first paragraph,
\(e^{-x_n}>1-x_n\), and \(x_n<1/(144n^2)<1/(12n^2)\), so
\(e^{-x_n}>1-(12n^2)^{-1}\).
\end{proof}
\end{annexitem}

\begin{annexitem}{annex:p15-w02c-p15-011-dml-from-clt}{thm:p15-w02c-p15-011-dml-from-clt}{Specialization to fair signs}
\begin{proof}
Each \(\varepsilon_i\) takes the values \(\pm1\) with probability \(1/2\),
so \(\mathbb E\varepsilon_1=\tfrac12-\tfrac12=0\), and
\(\varepsilon_1^2=1\) gives
\(\operatorname{Var}(\varepsilon_1)=\mathbb E\varepsilon_1^2=1\). The signs
are independent and identically distributed, so \cref{thm:p15-clt-iid}
applies with \(\mu=0\) and \(\sigma=1\) and gives
\(S_N/\sqrt N\Rightarrow\chi\).

Write \(Z_N=S_N/\sqrt N\) and \(F_N(x)=\mathbb P(Z_N\leq x)\). The function
\(\Phi\) is continuous at every real point: for \(x<y\),
\(0\leq\Phi(y)-\Phi(x)=\int_x^y\phi\leq(y-x)/\sqrt{2\pi}\), since
\(0<\phi\leq(2\pi)^{-1/2}\). By \cref{def:p15-weak-law-convergence},
\(F_N(x)\to\Phi(x)\) for every real \(x\).

For the strict inequality at a real \(x\), let \(\delta>0\). Monotonicity
of probability gives
\(F_N(x-\delta)\leq\mathbb P(Z_N<x)\leq F_N(x)\). Hence
\[
 \Phi(x-\delta)\leq\liminf_{N}\mathbb P(Z_N<x)
   \leq\limsup_{N}\mathbb P(Z_N<x)\leq\Phi(x),
\]
and letting \(\delta\downarrow0\), by continuity of \(\Phi\), gives
\(\mathbb P(Z_N<x)\to\Phi(x)\).

Now let \(a\leq b\) be real. The events are differences of nested events:
\[
 \begin{aligned}
 \mathbb P(a\leq Z_N\leq b)&=F_N(b)-\mathbb P(Z_N<a),\\
 \mathbb P(a<Z_N\leq b)&=F_N(b)-F_N(a),\\
 \mathbb P(a\leq Z_N<b)&=\mathbb P(Z_N<b)-\mathbb P(Z_N<a),\\
 \mathbb P(a<Z_N<b)&=\mathbb P(Z_N<b)-F_N(a)\qquad(a<b).
 \end{aligned}
\]
The first three hold for every \(a\leq b\). When \(a=b\) the event
\(\{a<Z_N<b\}\) is empty, with probability \(0=\Phi(b)-\Phi(a)\). By the two preceding paragraphs each right
side tends to \(\Phi(b)-\Phi(a)\). Since \(Z_N\) is real, an endpoint
\(-\infty\) imposes no condition, so the corresponding term
\(\mathbb P(Z_N<a)\) or \(F_N(a)\) is replaced by \(0=\Phi(-\infty)\); an
endpoint \(+\infty\) likewise replaces \(F_N(b)\) or \(\mathbb P(Z_N<b)\) by
\(1=\Phi(+\infty)\). If \(a=b=\pm\infty\), the event is empty and both sides
are zero. Finally \(\Phi(b)-\Phi(a)=\int_a^b\phi\) by the definition of
\(\Phi\) and additivity of the integral over \((-\infty,a]\) and \((a,b]\),
with the improper integrals at infinite endpoints taken as increasing limits
of nonnegative integrals and \(\int_{\mathbb R}\phi=1\).
\end{proof}
\end{annexitem}

\begin{annexitem}{annex:p15-w02c-p15-011-cdf-other-order}{thm:p15-w02c-p15-011-cdf-other-order}{Inner evaluation and Tonelli}
\begin{proof}
\emph{The half-plane is product measurable.} Its complement is
\[
 \mathbb R^2\setminus H_z=\{(x,y):x+y>z\}
   =\bigcup_{q\in\mathbb Q}\,(q,\infty)\times(z-q,\infty).
\]
If \(x>q\) and \(y>z-q\), then \(x+y>q+(z-q)=z\). Conversely, if
\(x+y>z\), then \(z-y<x\), and a rational \(q\) with \(z-y<q<x\) gives
\(x>q\) and \(y>z-q\). The right side is a countable union of measurable
rectangles, so \(\mathbb R^2\setminus H_z\), and with it \(H_z\), belongs
to \(\mathcal B(\mathbb R)\times\mathcal B(\mathbb R)\)
(\cref{def:p15-product-events}).

\emph{The inner integral at fixed \(x\).} Fix \(x\in\mathbb R\). For real
\(y\), the pair \((x,y)\) lies in \(H_z\) exactly when \(x+y\leq z\), that
is, exactly when \(y\leq z-x\). The section function
\(y\mapsto\mathbf1_{H_z}(x,y)\) is therefore the indicator of the Borel set
\((-\infty,z-x]\), and the integral of an indicator is the measure of its
set:
\[
 \int_{\mathbb R}\mathbf1_{H_z}(x,y)\,d\nu(y)
   =\nu((-\infty,z-x])=G(z-x),
\]
the last equality by the definition of the cumulative function \(G\) of
\(\nu\), evaluated at the point \(z-x\).

\emph{The Tonelli step.} Probability measures are finite, hence
sigma-finite, and \(f=\mathbf1_{H_z}\) is a product measurable function with
values in \(\{0,1\}\subseteq[0,\infty]\). Apply \cref{thm:p15-tonelli} with
\(\mu_1=\mu\) on the first coordinate, \(\mu_2=\nu\) on the second, and
\(\rho=\mu\times\nu\). Its first section integral is
\(T_1f(x)=\int\mathbf1_{H_z}(x,y)\,d\nu(y)=G(z-x)\) by the previous
paragraph. Tonelli's theorem says that \(T_1f\) is measurable, which is the
asserted Borel measurability of \(x\mapsto G(z-x)\), and that
\[
 \int_{\mathbb R}G(z-x)\,d\mu(x)=\int_{\mathbb R}T_1f\,d\mu
   =\int_{\mathbb R^2}\mathbf1_{H_z}\,d(\mu\times\nu)
   =(\mu\times\nu)(H_z).
\]

\emph{Identification with the convolution.} By
\cref{def:p15-convolution-law} with \(A=(-\infty,z]\),
\[
 \begin{aligned}
 (F*G)(z)&=(\mu*\nu)((-\infty,z])\\
   &=(\mu\times\nu)(\{(x,y):x+y\in(-\infty,z]\})
   =(\mu\times\nu)(H_z).
 \end{aligned}
\]
Combining the two displays gives
\((F*G)(z)=(\mu\times\nu)(H_z)=\int G(z-x)\,d\mu(x)\). The notation
\(\int G(z-x)\,dF(x)\) denotes the same integral against \(\mu\), by the
convention of \cref{def:p15-convolution-law}, so all four expressions in
the statement are equal.
\end{proof}
\end{annexitem}

\begin{annexitem}{annex:p15-w02c-p15-011-beta-rescaling}{thm:p15-w02c-p15-011-beta-rescaling}{The substitution \(y=zu\)}
\begin{proof}
\emph{The integrands as Borel functions on the line.} Define
\(h:\mathbb R\to[0,\infty)\) by
\[
 h(y)=\begin{cases}
  (z-y)^{\alpha-1}y^{\beta-1},&0<y<z,\\
  0,&\text{otherwise}.
 \end{cases}
\]
On the open interval \((0,z)\) both \(z-y\) and \(y\) are positive, so the
real powers are defined, positive and continuous there; \(h\) is therefore
a nonnegative Borel function. By the zero-endpoint convention of
\cref{def:p15-beta-integral}, the left side of the statement is
\(\int_{\mathbb R}h\,d\lambda\). Likewise put
\[
 k(u)=\begin{cases}
  (1-u)^{\alpha-1}u^{\beta-1},&0<u<1,\\
  0,&\text{otherwise},
 \end{cases}
\]
so that \(\int_{\mathbb R}k\,d\lambda=\int_0^1(1-u)^{\alpha-1}u^{\beta-1}\,
d\lambda(u)\). This equals \(B(\beta,\alpha)\) by
\cref{def:p15-beta-integral} with the parameters taken in the order
\(\beta,\alpha\), whose integrand is \(u^{\beta-1}(1-u)^{\alpha-1}\).

\emph{The substitution and its scaling factor.} Apply
\cref{thm:p15-affine-lebesgue} to the nonnegative Borel function \(h\) with
\(c=z\) and \(d=0\); since \(z>0\), \(|c|=z\). It gives
\[
 \int_{\mathbb R}h(zu)\,d\lambda(u)=\frac1z\int_{\mathbb R}h(y)\,d\lambda(y),
\]
that is,
\[
 \int_{\mathbb R}h\,d\lambda=z\int_{\mathbb R}h(zu)\,d\lambda(u).
\]
The factor \(z\) in front of the right-hand integral is the scaling factor
\(\lambda(zA)=z\lambda(A)\) of the same theorem, the factor that Durrett
records as \(dy=x\,du\) \cite[p.~51, (2.1.1)]{DurrettPTEAuthor2019}. Both
sides are nonnegative extended reals, and multiplication by the positive
number \(z\) is valid in \([0,\infty]\).

\emph{Rewriting the factors.} Because \(z>0\), the point \(zu\) lies in
\((0,z)\) exactly when \(u\in(0,1)\). For \(u\notin(0,1)\),
\(h(zu)=0=z^{\alpha-1}z^{\beta-1}k(u)\). For \(u\in(0,1)\),
\(z-zu=z(1-u)\) with \(z>0\) and \(1-u>0\), and \(zu\) has the positive
factors \(z\) and \(u\). For positive reals \(s,t\) and real \(p\),
\((st)^p=e^{p\log(st)}=e^{p\log s+p\log t}=s^pt^p\). Applying this with
\(p=\alpha-1\) and with \(p=\beta-1\),
\[
 h(zu)=(z-zu)^{\alpha-1}(zu)^{\beta-1}
      =z^{\alpha-1}(1-u)^{\alpha-1}\,z^{\beta-1}u^{\beta-1}
      =z^{\alpha-1}z^{\beta-1}\,k(u).
\]
So \(h(zu)=z^{\alpha-1}z^{\beta-1}k(u)\) for every real \(u\). A positive
constant factor comes out of a nonnegative integral, hence
\[
 \int_{\mathbb R}h(zu)\,d\lambda(u)
   =z^{\alpha-1}z^{\beta-1}\int_{\mathbb R}k\,d\lambda .
\]

\emph{The exponent bookkeeping.} For \(z>0\) and real \(p,q\),
\(z^pz^q=e^{p\log z}e^{q\log z}=e^{(p+q)\log z}=z^{p+q}\), and \(z=z^1\).
Hence
\[
 z\cdot z^{\alpha-1}\cdot z^{\beta-1}
   =z^{1+(\alpha-1)+(\beta-1)}=z^{\alpha+\beta-1}.
\]
Combining the three displays,
\[
 \begin{aligned}
 \int_0^z(z-y)^{\alpha-1}y^{\beta-1}\,d\lambda(y)
  &=\int_{\mathbb R}h\,d\lambda
   =z\cdot z^{\alpha-1}z^{\beta-1}\int_{\mathbb R}k\,d\lambda\\
  &=z^{\alpha+\beta-1}\int_0^1(1-u)^{\alpha-1}u^{\beta-1}\,d\lambda(u)\\
  &=z^{\alpha+\beta-1}B(\beta,\alpha).
 \end{aligned}
\]
No finiteness was used: every step is an identity in \([0,\infty]\).
\end{proof}
\end{annexitem}

\begin{annexitem}{annex:p15-w02c-p15-011-zero-variance-sum}{thm:p15-w02c-p15-011-zero-variance-sum}{An a.s.\ constant summand}
\begin{proof}
\emph{\(Y=\nu\) almost surely.} The law of \(Y\) is \(\delta_\nu\), so
\(\mathbb P(Y=\nu)=\delta_\nu(\{\nu\})=1\). Write
\(N=\{Y\neq\nu\}\); then \(\mathbb P(N)=0\).

\emph{\(X+Y=X+\nu\) almost surely.} On \(\Omega\setminus N\),
\(Y(\omega)=\nu\) and so \(X(\omega)+Y(\omega)=X(\omega)+\nu\). The two
variables can differ only on \(N\).

\emph{The law of \(X+\nu\).} This is the affine image \(cX+d\) with \(c=1\)
and \(d=\nu\). If \(a>0\), \(X\) has density \(\phi_{\mu,a}\), and
\cref{thm:p15-affine-lebesgue} with \(c=1\), \(d=\nu\) gives \(X+\nu\) the
density
\[
 \begin{aligned}
 t\longmapsto\phi_{\mu,a}(t-\nu)
  &=\frac1{\sqrt{2\pi a}}\exp\left(-\frac{(t-\nu-\mu)^2}{2a}\right)\\
  &=\phi_{\mu+\nu,a}(t),
 \end{aligned}
\]
so \(X+\nu\) has law \(\operatorname{normal}(\mu+\nu,a)\). If \(a=0\), the
law of \(X\) is \(\delta_\mu\), so \(X=\mu\) almost surely by the argument of
the first paragraph applied to \(X\); then \(X+\nu=\mu+\nu\) almost surely,
and \(X+\nu\) has law \(\delta_{\mu+\nu}=\operatorname{normal}(\mu+\nu,0)\)
by \cref{def:p15-normal-law}. In both cases this is the affine assertion of
\cref{thm:p15-normal-sums} at \(c=1\), \(d=\nu\), since \(c\mu+d=\mu+\nu\)
and \(c^2a=a\).

\emph{Almost sure equality preserves the law.} Let \(A\) be a Borel set.
The events \(\{X+Y\in A\}\) and \(\{X+\nu\in A\}\) agree outside \(N\), so
each is contained in the union of the other with \(N\). Monotonicity and
subadditivity give
\[
 \mathbb P(X+Y\in A)\leq\mathbb P(X+\nu\in A)+\mathbb P(N)
   =\mathbb P(X+\nu\in A),
\]
and the reverse inequality in the same way. Hence
\(\mathbb P(X+Y\in A)=\mathbb P(X+\nu\in A)\) for every Borel \(A\), and
\(X+Y\) has the law \(\operatorname{normal}(\mu+\nu,a)\) of \(X+\nu\). As
\(a=a+0\), this is \(\operatorname{normal}(\mu+\nu,a+b)\) at \(b=0\).
Independence of \(X\) and \(Y\) was not used.

\emph{\(Y\) is independent of \(X\).} By
\cref{def:p15-independent-random-elements} for the two variables, it
suffices to show
\(\mathbb P(X\in A,Y\in B)=\mathbb P(X\in A)\mathbb P(Y\in B)\) for Borel
\(A,B\); the conditions for subfamilies of at most one variable hold
trivially. If \(\nu\notin B\), then \(\{Y\in B\}\subseteq N\), so both sides
are zero. If \(\nu\in B\), then \(\mathbb P(Y\in B)=1\) and
\(\{X\in A\}\setminus\{X\in A,Y\in B\}\subseteq\{Y\notin B\}\subseteq N\),
so \(\mathbb P(X\in A,Y\in B)=\mathbb P(X\in A)
=\mathbb P(X\in A)\mathbb P(Y\in B)\).
\end{proof}
\end{annexitem}
